% Source: part_of_speech.tex
\begin{notechapter}
\chapter{The Viterbi Algorithm for Part-of-Speech Tagging}\label{ch:viterbi-pos}

\section{Introduction}
The Viterbi algorithm is a dynamic programming algorithm 
used to find the most likely sequence of hidden states 
(e.g., POS tags) given a sequence of observations (e.g., 
words) in a Hidden Markov Model (HMM).

\section{Key Components of an HMM}
\begin{itemize}
    \item Observations: Words in a sentence (e.g., "The cat 
          sat on the floor").
    \item Hidden States: POS tags (e.g., Det, Noun, Verb, 
          Prep).
    \item Transition Probabilities: \( P(s_i | s_{i-1}) \).
    \item Emission Probabilities: \( P(w_i | s_i) \).
    \item Initial Probabilities: \( P(s_1 | \text{start}) \).
\end{itemize}

\section{Probability Tables}
\subsection{Transition Probabilities}
\[
\begin{array}{|c|c|c|c|c|}
\hline
 & \text{Det} & \text{Noun} & \text{Verb} & \text{Prep} \\
\hline
\text{Start} & 0.5 & 0.3 & 0.1 & 0.1 \\
\text{Det} & 0.0 & 0.75 & 0.0 & 0.25 \\
\text{Noun} & 0.0 & 0.0 & 0.55 & 0.45 \\
\text{Verb} & 0.4 & 0.0 & 0.0 & 0.6 \\
\text{Prep} & 0.7 & 0.3 & 0.0 & 0.0 \\
\hline
\end{array}
\]

\subsection{Emission Probabilities}
\[
\begin{array}{|c|c|c|c|c|c|c|}
\hline
 & \text{"The"} & \text{"cat"} & \text{"sat"} & \text{"on"} & \text{"the"} & \text{"floor"} \\
\hline
\text{Det} & 0.6 & 0.0 & 0.0 & 0.0 & 0.4 & 0.0 \\
\text{Noun} & 0.0 & 0.5 & 0.0 & 0.0 & 0.0 & 0.4 \\
\text{Verb} & 0.0 & 0.0 & 0.7 & 0.0 & 0.0 & 0.3 \\
\text{Prep} & 0.0 & 0.0 & 0.0 & 0.8 & 0.0 & 0.0 \\
\hline
\end{array}
\]

\section{HMM Graph}
Below is the graph of the HMM showing transitions (blue arrows) and 
emissions (red arrows):

\begin{center}
\begin{tikzpicture}[
    > = Stealth, % Arrow head style
    shorten > = 1pt, % Don't touch arrow head to node
    auto,
    node distance = 3cm, % Distance between nodes
    semithick, % Line style
]

    % Nodes for states
    \node[state, color=green] (start) {Start};
    \node[state, color=blue] (Det) [below of=start, yshift=-1.5cm][left of=start, xshift=-1.5cm] {Det};
    \node[state, color=blue] (Noun) [right of=Det] {Noun};
    \node[state, color=blue] (Verb) [right of=Noun] {Verb};
    \node[state, color=blue] (Prep) [right of=Verb] {Prep};

    % Transitions
    \path[->, color=green] (start) edge node {0.5} (Det);
    \path[->, color=green] (start) edge node {0.3} (Noun);
    \path[->, color=green] (start) edge node {0.1} (Verb);
    \path[->, color=green] (start) edge node {0.1} (Prep);
    \path[->, color=blue] (Det) edge node {0.75} (Noun);
    \path[->, color=blue] (Det) edge [bend right] node {0.25} (Prep); 
    \path[->, color=blue] (Noun) edge node {0.55} (Verb); 
    \path[->, color=blue] (Noun) edge [bend right] node {0.45} (Prep);
    \path[->, color=blue] (Verb) edge [bend right] node {0.4} (Det);
    \path[->, color=blue] (Verb) edge node {0.6} (Prep);
    \path[->, color=blue] (Prep) edge [bend right] node {0.7} (Det);
    \path[->, color=blue] (Prep) edge [bend right] node {0.3} (Noun);

    % Emissions
    \node[state, color=red] (The) [below of=Det, yshift=-1cm][left of=Det] {"The"};
    \path[->, color=red] (Det) edge node {0.6} (The);
    \node[state, color=red] (the) [below of=Det, yshift=-1cm] {"the"};
    \path[->, color=red] (Det) edge node {0.4} (the);
    \node[state, color=red] (cat) [below of=Noun, xshift=-1cm, yshift=-1cm] {"cat"};
    \path[->, color=red] (Noun) edge node {0.5} (cat);
    \node[state, color=red] (sat) [below of=Verb, xshift=1cm, yshift=-1cm] {"sat"};
    \path[->, color=red] (Verb) edge node {0.7} (sat);
    \node[state, color=red] (on) [below of=Prep, yshift=-1cm] {"on"};
    \path[->, color=red] (Prep) edge node {0.8} (on);
    \node[state, color=red] (floor) [below of=Noun, xshift=2cm, yshift=-1cm] {"floor"};
    \path[->, color=red] (Noun) edge node {0.4} (floor);
    \path[->, color=red] (Verb) edge node {0.3} (floor);
\end{tikzpicture}
\end{center}

\section{Why the Viterbi Algorithm Works}
The Viterbi algorithm works because it leverages the 
\textbf{optimal substructure} and \textbf{overlapping subproblems} 
properties of dynamic programming. Here's why:

\subsection{Optimal Substructure}
The problem of finding the most likely sequence of hidden 
states can be broken down into smaller subproblems. For 
example, the most likely path to a state \( s \) at time \( t \) 
depends only on:
\begin{itemize}
    \item The most likely path to each possible state at 
          time \( t-1 \).
    \item The transition probability from those states to 
          \( s \).
    \item The emission probability of the observation at 
          time \( t \) given \( s \).
\end{itemize}

This means that the optimal solution to the problem can be 
constructed from optimal solutions to its subproblems.

\subsection{Overlapping Subproblems}
The Viterbi algorithm avoids redundant computations by 
storing intermediate results in a table (the trellis). For 
example, the probability \( \delta_t(s) \) is computed once 
and reused for all future steps. This makes the algorithm 
efficient, with a time complexity of \( O(T \cdot N^2) \), 
where \( T \) is the number of observations and \( N \) is the 
number of states.

\subsection{Backtracking}
Once the trellis is filled, the algorithm backtracks to 
reconstruct the most likely sequence of states. This is 
possible because the trellis stores not only the maximum 
probabilities but also the states that led to them.

\section{Complete Trellis with Arrows and Probabilities}
The trellis is a table where rows represent states (tags) 
and columns represent words. Each cell stores the 
probability of the most likely path to that state at that 
step.

\begin{center}
\begin{tikzpicture}[
    node distance=1.3cm,
    every node/.style={rectangle, draw, minimum width=1.2cm, 
    minimum height=0.9cm, align=center},
    arrow/.style={->, >=Stealth, thick}
]

    % Words (Observations)
    \node (The) {The};
    \node (cat) [right=of The] {cat};
    \node (sat) [right=of cat] {sat};
    \node (on) [right=of sat] {on};
    \node (the) [right=of on] {the};
    \node (floor) [right=of the] {floor};

    % States (Tags)
    \node (Det1) [below=of The] {Det \\ 0.3};
    \node (Noun1) [below=of Det1] {Noun \\ 0.0};
    \node (Verb1) [below=of Noun1] {Verb \\ 0.0};
    \node (Prep1) [below=of Verb1] {Prep \\ 0.0};

    \node (Det2) [below=of cat] {Det \\ 0.0};
    \node (Noun2) [below=of Det2] {Noun \\ 0.1125};
    \node (Verb2) [below=of Noun2] {Verb \\ 0.0};
    \node (Prep2) [below=of Verb2] {Prep \\ 0.0};

    \node (Det3) [below=of sat] {Det \\ 0.0};
    \node (Noun3) [below=of Det3] {Noun \\ 0.0};
    \node (Verb3) [below=of Noun3] {Verb \\ 0.0433};
    \node (Prep3) [below=of Verb3] {Prep \\ 0.0};

    \node (Det4) [below=of on] {Det \\ 0.0};
    \node (Noun4) [below=of Det4] {Noun \\ 0.0};
    \node (Verb4) [below=of Noun4] {Verb \\ 0.0};
    \node (Prep4) [below=of Verb4] {Prep \\ 0.02079};

    \node (Det5) [below=of the] {Det \\ 0.005821};
    \node (Noun5) [below=of Det5] {Noun \\ 0.0};
    \node (Verb5) [below=of Noun5] {Verb \\ 0.0};
    \node (Prep5) [below=of Verb5] {Prep \\ 0.0};

    \node (Det6) [below=of floor] {Det \\ 0.0};
    \node (Noun6) [below=of Det6] {Noun \\ 0.0017};
    \node (Verb6) [below=of Noun6] {Verb \\ 0.0};
    \node (Prep6) [below=of Verb6] {Prep \\ 0.0};

    % Arrows for paths
    \draw[arrow] (Det1) -- (Noun2);
    \draw[arrow] (Noun2) -- (Verb3);
    \draw[arrow] (Verb3) -- (Prep4);
    \draw[arrow] (Prep4) -- (Det5);
    \draw[arrow] (Det5) -- (Noun6);

    % Labels for states
    \node[left=of Det1] {Det};
    \node[left=of Noun1] {Noun};
    \node[left=of Verb1] {Verb};
    \node[left=of Prep1] {Prep};
\end{tikzpicture}
\end{center}

\section{Step-by-Step Computations}
\subsection{General Formula for \(\delta_t(s)\)}
\[
\delta_t(s) = \max_{s'} \left( \delta_{t-1}(s') \cdot P(s | s') \cdot P(w_t | s) \right)
\]
where:
\begin{itemize}
    \item \(\delta_t(s)\): Probability of the most likely path to state \(s\) at time \(t\).
    \item \(P(s | s')\): Transition probability from state \(s'\) to \(s\).
    \item \(P(w_t | s)\): Emission probability of word \(w_t\) given state \(s\).
\end{itemize}

\subsubsection{Step 1: Initialization}
For \( w_1 = \text{"The"} \):
\[
\delta_1(\text{Det}) = 0.5 \cdot 0.6 = 0.3
\]
\[
\delta_1(\text{Noun}) = 0.3 \cdot 0.0 = 0.0
\]
\[
\delta_1(\text{Verb}) = 0.1 \cdot 0.0 = 0.0
\]
\[
\delta_1(\text{Prep}) = 0.1 \cdot 0.0 = 0.0
\]

\subsubsection{Step 2: Recursion for \( w_2 = \text{"cat"} \)}
For \( s = \text{Noun} \):
\[
\delta_2(\text{Noun}) = 0.3 \cdot 0.75 \cdot 0.5 = 0.1125
\]

\subsubsection{Step 3: Recursion for \( w_3 = \text{"sat"} \)}
For \( s = \text{Verb} \):
\[
\delta_3(\text{Verb}) = 0.1125 \cdot 0.55 \cdot 0.7 = 0.0433
\]

\subsubsection{Step 4: Recursion for \( w_4 = \text{"on"} \)}
For \( s = \text{Prep} \):
\[
\delta_4(\text{Prep}) = 0.0433 \cdot 0.6 \cdot 0.8 = 0.02079
\]

\subsubsection{Step 5: Recursion for \( w_5 = \text{"the"} \)}
For \( s = \text{Det} \):
\[
\delta_5(\text{Det}) = 0.02079 \cdot 0.7 \cdot 0.4 = 0.005821
\]

\subsubsection{Step 6: Recursion for \( w_6 = \text{"floor"} \)}
For \( s = \text{Noun} \):
\[
\delta_6(\text{Noun}) = 0.005821 \cdot 0.75 \cdot 0.4 \approx 0.001746
\]

\subsubsection{Step 7: Termination}
The best path probability is \( \max(0.001746, 0.0) \approx 0.001746 \) (\( \delta_6(\text{Verb}) = 0 \) since \( P(\text{Verb} \mid \text{Det}) = 0 \)).

\subsubsection{Step 8: Backtracking}
Trace back the most likely sequence:
\[
\text{Det} \rightarrow \text{Noun} \rightarrow \text{Verb} 
\rightarrow \text{Prep} \rightarrow \text{Det} \rightarrow 
\text{Noun}
\]

\section{Conclusion}
The Viterbi algorithm efficiently computes the most likely 
sequence of hidden states using dynamic programming and 
backtracking. The HMM graph and trellis visualization 
clearly show the transitions and the most likely path.
\end{notechapter}
